% Options for packages loaded elsewhere
\PassOptionsToPackage{unicode}{hyperref}
\PassOptionsToPackage{hyphens}{url}
%
\documentclass[
]{article}
\usepackage{amsmath,amssymb}
\usepackage{lmodern}
\usepackage{iftex}
\ifPDFTeX
  \usepackage[T1]{fontenc}
  \usepackage[utf8]{inputenc}
  \usepackage{textcomp} % provide euro and other symbols
\else % if luatex or xetex
  \usepackage{unicode-math}
  \defaultfontfeatures{Scale=MatchLowercase}
  \defaultfontfeatures[\rmfamily]{Ligatures=TeX,Scale=1}
\fi
% Use upquote if available, for straight quotes in verbatim environments
\IfFileExists{upquote.sty}{\usepackage{upquote}}{}
\IfFileExists{microtype.sty}{% use microtype if available
  \usepackage[]{microtype}
  \UseMicrotypeSet[protrusion]{basicmath} % disable protrusion for tt fonts
}{}
\makeatletter
\@ifundefined{KOMAClassName}{% if non-KOMA class
  \IfFileExists{parskip.sty}{%
    \usepackage{parskip}
  }{% else
    \setlength{\parindent}{0pt}
    \setlength{\parskip}{6pt plus 2pt minus 1pt}}
}{% if KOMA class
  \KOMAoptions{parskip=half}}
\makeatother
\usepackage{xcolor}
\IfFileExists{xurl.sty}{\usepackage{xurl}}{} % add URL line breaks if available
\IfFileExists{bookmark.sty}{\usepackage{bookmark}}{\usepackage{hyperref}}
\hypersetup{
  hidelinks,
  pdfcreator={LaTeX via pandoc}}
\urlstyle{same} % disable monospaced font for URLs
\usepackage{longtable,booktabs,array}
\usepackage{calc} % for calculating minipage widths
% Correct order of tables after \paragraph or \subparagraph
\usepackage{etoolbox}
\makeatletter
\patchcmd\longtable{\par}{\if@noskipsec\mbox{}\fi\par}{}{}
\makeatother
% Allow footnotes in longtable head/foot
\IfFileExists{footnotehyper.sty}{\usepackage{footnotehyper}}{\usepackage{footnote}}
\makesavenoteenv{longtable}
\setlength{\emergencystretch}{3em} % prevent overfull lines
\providecommand{\tightlist}{%
  \setlength{\itemsep}{0pt}\setlength{\parskip}{0pt}}
\setcounter{secnumdepth}{-\maxdimen} % remove section numbering
\ifLuaTeX
  \usepackage{selnolig}  % disable illegal ligatures
\fi

\author{}
\date{}

\begin{document}

\hypertarget{rl-medical-diagnosis-complete-project-explanation}{%
\section{🩺 RL Medical Diagnosis --- Complete Project
Explanation}\label{rl-medical-diagnosis-complete-project-explanation}}

\begin{quote}
\textbf{Read this document carefully before your viva/presentation.}
Every concept, every algorithm, every design choice is explained here in
simple English. After reading this, you will be able to answer ANY
question about your project confidently.
\end{quote}

\begin{center}\rule{0.5\linewidth}{0.5pt}\end{center}

\hypertarget{table-of-contents}{%
\subsection{Table of Contents}\label{table-of-contents}}

\begin{enumerate}
\def\labelenumi{\arabic{enumi}.}
\tightlist
\item
  \protect\hyperlink{1-the-problem-statement}{The Problem Statement ---
  What Are We Actually Doing?}
\item
  \protect\hyperlink{2-real-world-significance}{Real-World Significance
  --- Why This Matters}
\item
  \protect\hyperlink{3-what-is-reinforcement-learning}{What is
  Reinforcement Learning? (The Basics)}
\item
  \protect\hyperlink{4-mdp-formulation}{MDP Formulation --- The Heart of
  the Project}
\item
  \protect\hyperlink{5-transition-probability}{Transition Probability
  --- The Most Important Calculation}
\item
  \protect\hyperlink{6-the-10-algorithms}{The 10 Algorithms ---
  Explained One by One}
\item
  \protect\hyperlink{7-why-policy-iteration-wins}{Why Policy Iteration
  Wins Over Everything Else}
\item
  \protect\hyperlink{8-results-discussion}{Results Discussion}
\item
  \protect\hyperlink{9-likely-viva-questions-and-answers}{Likely Viva
  Questions and Answers}
\end{enumerate}

\begin{center}\rule{0.5\linewidth}{0.5pt}\end{center}

\hypertarget{the-problem-statement}{%
\subsection{1. The Problem Statement}\label{the-problem-statement}}

\hypertarget{in-one-sentence}{%
\subsubsection{In One Sentence}\label{in-one-sentence}}

\begin{quote}
We built an AI doctor that learns to diagnose diseases by asking the
right questions in the right order.
\end{quote}

\hypertarget{the-full-picture}{%
\subsubsection{The Full Picture}\label{the-full-picture}}

Imagine you walk into a doctor's office. The doctor doesn't know what's
wrong with you. They don't run every test in the world --- that would be
expensive and time-consuming. Instead, a good doctor asks \emph{smart
questions}:

\begin{itemize}
\tightlist
\item
  ``Do you have a fever?'' → If yes, many diseases eliminated. If no,
  different set eliminated.
\item
  Based on that answer, they ask the \emph{next best question}.
\item
  After 3-4 good questions, they can confidently tell you what you have.
\end{itemize}

\textbf{Our project does exactly this, but with AI.} We trained an AI
agent to learn this skill using Reinforcement Learning.

\hypertarget{the-setup}{%
\subsubsection{The Setup}\label{the-setup}}

\begin{itemize}
\tightlist
\item
  \textbf{8 diseases} the AI can diagnose: Flu, Strep, Pneumonia,
  Bronchitis, Cold, Allergy, Asthma, Migraine
\item
  \textbf{5 symptoms} the AI can ask about: Fever, Cough, Fatigue,
  Breathing Difficulty, Headache
\item
  Each disease has a \textbf{unique combination} of symptoms (see table
  below)
\end{itemize}

\begin{longtable}[]{@{}llllll@{}}
\toprule
Disease & Fever & Cough & Fatigue & Breath & Headache \\
\midrule
\endhead
Flu & ✓ & ✓ & ✓ & & \\
Strep & ✓ & ✓ & & ✓ & \\
Pneumonia & ✓ & & ✓ & & ✓ \\
Bronchitis & ✓ & & & ✓ & ✓ \\
Cold & & ✓ & ✓ & & ✓ \\
Allergy & & ✓ & & ✓ & \\
Asthma & & & ✓ & ✓ & \\
Migraine & & & & & ✓ \\
\bottomrule
\end{longtable}

\hypertarget{the-challenge}{%
\subsubsection{The Challenge}\label{the-challenge}}

Look at that table carefully. Flu and Strep both have Fever AND Cough.
If you only know ``Fever=Yes, Cough=Yes'', you still can't tell if it's
Flu or Strep! You need to ask at least one more question (Fatigue or
Breath) to distinguish them.

\textbf{The AI must learn:} 1. Which symptom to ask FIRST (the one that
eliminates the most diseases) 2. Which symptom to ask NEXT based on what
it learned 3. WHEN to stop asking and commit to a diagnosis 4. All of
this while minimizing the number of questions (because each question has
a small cost)

\hypertarget{why-not-just-ask-all-5-questions-every-time}{%
\subsubsection{Why Not Just Ask All 5 Questions Every
Time?}\label{why-not-just-ask-all-5-questions-every-time}}

Because in the real world, every test costs money and time. If the AI
can diagnose correctly after only 3-4 questions instead of 5, that saves
the patient money, time, and unnecessary tests. The AI needs to balance
\textbf{accuracy vs.~efficiency}.

\begin{center}\rule{0.5\linewidth}{0.5pt}\end{center}

\hypertarget{real-world-significance}{%
\subsection{2. Real-World Significance}\label{real-world-significance}}

\hypertarget{where-would-this-actually-be-used}{%
\subsubsection{Where Would This Actually Be
Used?}\label{where-would-this-actually-be-used}}

\begin{enumerate}
\def\labelenumi{\arabic{enumi}.}
\item
  \textbf{Hospital Triage}: When a patient walks into an emergency room,
  a nurse needs to quickly figure out how serious their condition is.
  Our AI could help prioritize --- ``Ask about fever first, then
  breathing difficulty'' --- and guide the triage process.
\item
  \textbf{Rural Healthcare}: In remote areas with limited access to
  specialists, a community health worker with this AI tool on a
  phone/tablet could do preliminary diagnosis by asking the right
  questions.
\item
  \textbf{Reducing Unnecessary Tests}: In real medicine, blood tests,
  X-rays, and MRIs are expensive. If the AI determines that just asking
  3 questions is enough to diagnose, it saves the patient from 2
  unnecessary tests.
\item
  \textbf{Clinical Decision Support}: Doctors are human --- they can
  forget, get tired, or have biases. An AI that always follows the
  optimal question order provides a consistent second opinion.
\item
  \textbf{Beyond Healthcare}: The same idea applies to IT tech support
  (``Is your computer turning on?'' → ``Is the screen showing
  anything?''), customer service chatbots, and quality control in
  factories.
\end{enumerate}

\hypertarget{why-reinforcement-learning-and-not-simple-rules}{%
\subsubsection{Why Reinforcement Learning and Not Simple
Rules?}\label{why-reinforcement-learning-and-not-simple-rules}}

You might ask: \emph{``Why not just hard-code a decision tree? If fever,
then ask cough. If cough, then\ldots{}''}

Three reasons: 1. \textbf{Adaptability}: If you add new diseases or new
symptoms, you'd have to rewrite all the rules. RL just retrains
automatically. 2. \textbf{Optimality}: RL mathematically guarantees
finding the BEST question order. A human-designed decision tree might
miss a more efficient ordering. 3. \textbf{Balancing Trade-offs}: RL
naturally handles the cost-accuracy trade-off through its reward
structure. Hard-coding this is very difficult.

\begin{center}\rule{0.5\linewidth}{0.5pt}\end{center}

\hypertarget{what-is-reinforcement-learning}{%
\subsection{3. What is Reinforcement
Learning?}\label{what-is-reinforcement-learning}}

\hypertarget{the-core-idea}{%
\subsubsection{The Core Idea}\label{the-core-idea}}

Think of training a dog: - The dog (agent) tries doing something - If it
does something good, you give it a treat (positive reward) - If it does
something bad, you say ``No!'' (negative reward) - Over time, the dog
learns what actions lead to treats

\textbf{RL is the same thing, but for computers.}

\hypertarget{the-key-components}{%
\subsubsection{The Key Components}\label{the-key-components}}

\begin{longtable}[]{@{}
  >{\raggedright\arraybackslash}p{(\columnwidth - 4\tabcolsep) * \real{0.28}}
  >{\raggedright\arraybackslash}p{(\columnwidth - 4\tabcolsep) * \real{0.32}}
  >{\raggedright\arraybackslash}p{(\columnwidth - 4\tabcolsep) * \real{0.40}}@{}}
\toprule
Component & Dog Training & Our Medical AI \\
\midrule
\endhead
\textbf{Agent} & The dog & The AI doctor \\
\textbf{Environment} & The room/park & The patient (with their hidden
disease) \\
\textbf{State} & What the dog sees/knows & What symptoms the AI has
asked about so far \\
\textbf{Action} & Sit, roll over, bark & Ask about Fever, Diagnose Flu,
etc. \\
\textbf{Reward} & Treat or ``No!'' & +10 for correct diagnosis, -5 for
wrong, -0.1 per question \\
\textbf{Policy} & The dog's learned behavior & The AI's strategy: ``In
this situation, do this action'' \\
\bottomrule
\end{longtable}

\hypertarget{the-goal}{%
\subsubsection{The Goal}\label{the-goal}}

Find the \textbf{optimal policy} --- the strategy that maximizes the
total reward over time. In our case, that means: ask the fewest
questions possible while always diagnosing correctly.

\begin{center}\rule{0.5\linewidth}{0.5pt}\end{center}

\hypertarget{mdp-formulation}{%
\subsection{4. MDP Formulation}\label{mdp-formulation}}

MDP stands for \textbf{Markov Decision Process}. It's the mathematical
framework that describes our problem precisely. Let's go through each
component.

\hypertarget{the-agent}{%
\subsubsection{4.1 The Agent}\label{the-agent}}

The \textbf{agent} is the AI doctor. It makes decisions (which symptom
to ask about, when to diagnose).

\hypertarget{the-environment}{%
\subsubsection{4.2 The Environment}\label{the-environment}}

The \textbf{environment} is the patient. The patient has a hidden
disease, and when the AI asks ``Do you have fever?'', the environment
(patient) answers honestly.

\hypertarget{states-what-the-ai-knows}{%
\subsubsection{4.3 States --- What the AI
Knows}\label{states-what-the-ai-knows}}

A \textbf{state} represents what the AI currently knows about the
patient.

Since there are 5 symptoms, and each symptom can be in one of 3
conditions: - \textbf{0 = Unknown} --- haven't asked yet - \textbf{1 =
Absent} --- asked, and the patient said ``No'' - \textbf{2 = Present}
--- asked, and the patient said ``Yes''

The total number of possible knowledge combinations = 3 × 3 × 3 × 3 × 3
= \textbf{3⁵ = 243 states}

Plus 1 \textbf{terminal state} (state 243) = the episode is over
(diagnosis has been made).

\textbf{Total: 244 states.}

\hypertarget{how-states-are-encoded-base-3}{%
\subsubsection{How States Are Encoded
(Base-3)}\label{how-states-are-encoded-base-3}}

We use a clever encoding. Each state is a single number calculated from
the 5 symptom statuses:

\begin{verbatim}
State = status[0] × 3⁰ + status[1] × 3¹ + status[2] × 3² + status[3] × 3³ + status[4] × 3⁴
\end{verbatim}

Where status{[}0{]} = Fever status, status{[}1{]} = Cough status, etc.

\textbf{Examples:}

\begin{longtable}[]{@{}lllllll@{}}
\toprule
State & Fever & Cough & Fatigue & Breath & Headache & Meaning \\
\midrule
\endhead
0 & 0 & 0 & 0 & 0 & 0 & Nothing known yet (start) \\
2 & 2 & 0 & 0 & 0 & 0 & Fever = Present, rest unknown \\
1 & 1 & 0 & 0 & 0 & 0 & Fever = Absent, rest unknown \\
8 & 2 & 2 & 0 & 0 & 0 & Fever = Present, Cough = Present \\
243 & --- & --- & --- & --- & --- & Terminal (episode done) \\
\bottomrule
\end{longtable}

For State 8: Fever=2, Cough=2, rest=0 → 2×1 + 2×3 = 2+6 = \textbf{8} ✓

\hypertarget{actions-what-the-ai-can-do}{%
\subsubsection{4.4 Actions --- What the AI Can
Do}\label{actions-what-the-ai-can-do}}

The AI has \textbf{13 possible actions}:

\begin{longtable}[]{@{}lll@{}}
\toprule
Action \# & Type & What It Does \\
\midrule
\endhead
0 & Ask & Ask about Fever \\
1 & Ask & Ask about Cough \\
2 & Ask & Ask about Fatigue \\
3 & Ask & Ask about Breath \\
4 & Ask & Ask about Headache \\
5 & Diagnose & Diagnose Flu \\
6 & Diagnose & Diagnose Strep \\
7 & Diagnose & Diagnose Pneumonia \\
8 & Diagnose & Diagnose Bronchitis \\
9 & Diagnose & Diagnose Cold \\
10 & Diagnose & Diagnose Allergy \\
11 & Diagnose & Diagnose Asthma \\
12 & Diagnose & Diagnose Migraine \\
\bottomrule
\end{longtable}

\textbf{The key design choice}: Having separate ``ask'' and ``diagnose''
actions forces the AI to learn TWO things: 1. WHAT to ask next 2. WHEN
to stop asking and make a diagnosis

This is the core trade-off: every question costs -0.1 reward, but
diagnosing wrong costs -5.0. So the AI must learn to ask enough
questions to be confident, but not more than necessary.

\hypertarget{rewards-the-ais-report-card}{%
\subsubsection{4.5 Rewards --- The AI's Report
Card}\label{rewards-the-ais-report-card}}

\begin{longtable}[]{@{}lll@{}}
\toprule
Event & Reward & Why \\
\midrule
\endhead
Ask any symptom & \textbf{-0.1} & Small penalty to discourage asking too
many questions \\
Correct diagnosis & \textbf{+10.0} & Big reward for getting it right \\
Wrong diagnosis & \textbf{-5.0} & Big punishment for getting it wrong \\
\bottomrule
\end{longtable}

\textbf{The discount factor γ = 0.9} means: future rewards are worth
90\% of immediate rewards. This makes the AI care about both the
immediate cost of asking and the future benefit of a correct diagnosis.

\textbf{Example}: If the AI asks 4 questions (cost = 4 × -0.1 = -0.4)
and then diagnoses correctly (+10.0), the total undiscounted return
would be about -0.4 + 10.0 = \textbf{+9.6}. If it asks all 5 questions
and diagnoses correctly: -0.5 + 10.0 = \textbf{+9.5} --- slightly less.
So fewer questions = better total reward, as long as accuracy is
maintained.

\hypertarget{policy-the-ais-strategy}{%
\subsubsection{4.6 Policy --- The AI's
Strategy}\label{policy-the-ais-strategy}}

A \textbf{policy} π(s) tells the AI: ``When you are in state s, take
action a.''

For example, a good policy might say: - State 0 (know nothing) → Action
0 (Ask Fever) - State 2 (Fever=Present) → Action 1 (Ask Cough) - State 8
(Fever=Present, Cough=Present) → Action 2 (Ask Fatigue) - State 17
(Fever=Present, Cough=Present, Fatigue=Absent) → Action 6 (Diagnose
Strep!)

The goal of all our algorithms is to find the \textbf{optimal policy}
--- the one that gives the maximum total reward.

\begin{center}\rule{0.5\linewidth}{0.5pt}\end{center}

\hypertarget{transition-probability-the-most-important-calculation}{%
\subsection{5. Transition Probability --- The Most Important
Calculation}\label{transition-probability-the-most-important-calculation}}

\begin{quote}
\textbf{This is what most people find confusing. Read this section very
carefully.}
\end{quote}

\hypertarget{what-is-a-transition-probability}{%
\subsubsection{What Is a Transition
Probability?}\label{what-is-a-transition-probability}}

P(s'\textbar s, a) means: ``If I am in state s and I take action a, what
is the probability that I end up in state s'?''

\hypertarget{why-do-we-need-it}{%
\subsubsection{Why Do We Need It?}\label{why-do-we-need-it}}

Dynamic Programming methods (Policy Iteration and Value Iteration) need
to plan ahead. They need to think: \emph{``If I ask about Fever right
now, what might happen? What state will I be in next?''}

But here's the catch: \textbf{the AI doesn't know which disease the
patient has!} So when it asks ``Do you have Fever?'', it doesn't know if
the answer will be Yes or No.~It depends on the hidden disease.

\hypertarget{how-we-calculate-it-step-by-step}{%
\subsubsection{How We Calculate It --- Step by
Step}\label{how-we-calculate-it-step-by-step}}

The key idea: \textbf{we average over all diseases that are still
possible.}

\hypertarget{case-1-starting-state-know-nothing}{%
\paragraph{CASE 1: Starting State (Know
Nothing)}\label{case-1-starting-state-know-nothing}}

\textbf{State = 0} (all symptoms unknown), \textbf{Action = ``Ask
Fever'' (a=0)}

Step 1: Which diseases are compatible? ALL 8 (since we know nothing).

Step 2: What happens if the patient has each disease?

\begin{longtable}[]{@{}lll@{}}
\toprule
Disease & Has Fever? & Next State \\
\midrule
\endhead
Flu & Yes (1) & Status becomes {[}2,0,0,0,0{]} → state \textbf{2} \\
Strep & Yes (1) & Status becomes {[}2,0,0,0,0{]} → state \textbf{2} \\
Pneumonia & Yes (1) & Status becomes {[}2,0,0,0,0{]} → state
\textbf{2} \\
Bronchitis & Yes (1) & Status becomes {[}2,0,0,0,0{]} → state
\textbf{2} \\
Cold & No (0) & Status becomes {[}1,0,0,0,0{]} → state \textbf{1} \\
Allergy & No (0) & Status becomes {[}1,0,0,0,0{]} → state \textbf{1} \\
Asthma & No (0) & Status becomes {[}1,0,0,0,0{]} → state \textbf{1} \\
Migraine & No (0) & Status becomes {[}1,0,0,0,0{]} → state \textbf{1} \\
\bottomrule
\end{longtable}

Step 3: Count the probabilities (each disease equally likely = 1/8):

\begin{itemize}
\tightlist
\item
  P(state 2 \textbar{} state 0, Ask Fever) = 4/8 = \textbf{0.5} (Flu,
  Strep, Pneumonia, Bronchitis have Fever)
\item
  P(state 1 \textbar{} state 0, Ask Fever) = 4/8 = \textbf{0.5} (Cold,
  Allergy, Asthma, Migraine don't)
\end{itemize}

\textbf{In plain English}: When you start and ask about Fever, there's a
50-50 chance the answer is Yes vs No.

\hypertarget{case-2-after-knowing-fever-present}{%
\paragraph{CASE 2: After Knowing Fever =
Present}\label{case-2-after-knowing-fever-present}}

\textbf{State = 2} (Fever=Present, rest unknown), \textbf{Action = ``Ask
Cough'' (a=1)}

Step 1: Compatible diseases = only those WITH Fever = \{Flu, Strep,
Pneumonia, Bronchitis\} (4 diseases).

Step 2: What happens for each?

\begin{longtable}[]{@{}lll@{}}
\toprule
Disease & Has Cough? & Next State \\
\midrule
\endhead
Flu & Yes (1) & Status {[}2,2,0,0,0{]} → state \textbf{8} \\
Strep & Yes (1) & Status {[}2,2,0,0,0{]} → state \textbf{8} \\
Pneumonia & No (0) & Status {[}2,1,0,0,0{]} → state \textbf{5} \\
Bronchitis & No (0) & Status {[}2,1,0,0,0{]} → state \textbf{5} \\
\bottomrule
\end{longtable}

Step 3: Each is equally likely (1/4 each):

\begin{itemize}
\tightlist
\item
  P(state 8 \textbar{} state 2, Ask Cough) = 2/4 = \textbf{0.5}
\item
  P(state 5 \textbar{} state 2, Ask Cough) = 2/4 = \textbf{0.5}
\end{itemize}

\hypertarget{case-3-diagnose-action}{%
\paragraph{CASE 3: Diagnose Action}\label{case-3-diagnose-action}}

\textbf{Any state, Action = ``Diagnose Flu'' (a=5)}

\begin{itemize}
\tightlist
\item
  The transition is always to \textbf{state 243} (terminal) with
  probability \textbf{1.0}
\item
  The REWARD depends on whether the diagnosis is correct:

  \begin{itemize}
  \tightlist
  \item
    If the patient actually has Flu → reward = +10.0
  \item
    If the patient has something else → reward = -5.0
  \end{itemize}
\end{itemize}

\hypertarget{case-4-after-knowing-feveryes-and-coughyes}{%
\paragraph{CASE 4: After Knowing Fever=Yes AND
Cough=Yes}\label{case-4-after-knowing-feveryes-and-coughyes}}

\textbf{State = 8} (Fever=Present, Cough=Present), \textbf{Action =
``Ask Fatigue'' (a=2)}

Step 1: Compatible = only diseases with Fever=Yes AND Cough=Yes = \{Flu,
Strep\} (2 diseases)

Step 2: \textbar{} Disease \textbar{} Has Fatigue? \textbar{} Next State
\textbar{}
\textbar---------\textbar-------------\textbar------------\textbar{}
\textbar{} Flu \textbar{} Yes (1) \textbar{} {[}2,2,2,0,0{]} → state
\textbf{26} \textbar{} \textbar{} Strep \textbar{} No (0) \textbar{}
{[}2,2,1,0,0{]} → state \textbf{17} \textbar{}

Step 3: - P(state 26 \textbar{} state 8, Ask Fatigue) = 1/2 =
\textbf{0.5} - P(state 17 \textbar{} state 8, Ask Fatigue) = 1/2 =
\textbf{0.5}

Now if we are in state 26 (Fever=Yes, Cough=Yes, Fatigue=Yes), only Flu
matches → \textbf{Diagnose Flu with confidence!} If we are in state 17
(Fever=Yes, Cough=Yes, Fatigue=No), only Strep matches →
\textbf{Diagnose Strep with confidence!}

\hypertarget{case-5-deep-decision-distinguishing-pneumonia-from-bronchitis}{%
\paragraph{CASE 5: Deep Decision --- Distinguishing Pneumonia from
Bronchitis}\label{case-5-deep-decision-distinguishing-pneumonia-from-bronchitis}}

Suppose we know: Fever=Yes (state 2) and then Cough=No (new state 5).

Compatible diseases = Fever=Yes AND Cough=No = \{Pneumonia, Bronchitis\}
(2 diseases)

\textbf{Action = ``Ask Fatigue'' (a=2)}: \textbar{} Disease \textbar{}
Has Fatigue? \textbar{} Next State \textbar{}
\textbar---------\textbar-------------\textbar------------\textbar{}
\textbar{} Pneumonia \textbar{} Yes (1) \textbar{} {[}2,1,2,0,0{]} →
different state \textbar{} \textbar{} Bronchitis \textbar{} No (0)
\textbar{} {[}2,1,1,0,0{]} → different state \textbar{}

Result: P = 0.5 for each outcome. - If Fatigue=Yes → only Pneumonia
compatible → Diagnose Pneumonia! - If Fatigue=No → only Bronchitis
compatible → Diagnose Bronchitis!

\hypertarget{the-general-formula}{%
\subsubsection{The General Formula}\label{the-general-formula}}

\begin{verbatim}
P(s' | s, a) = (Number of compatible diseases that lead to s') / (Total compatible diseases)
\end{verbatim}

This is the formula from the presentation:

P(s'\textbar s,a) = Σ {[}1/\textbar Compatible\textbar{]} × I{[}s' =
T(s,a,d){]} for all diseases d in Compatible

Where: - Compatible = diseases that match all known symptoms so far -
T(s,a,d) = the next state if disease is d and we take action a -
I{[}\ldots{]} = 1 if the condition is true, 0 otherwise

\hypertarget{why-this-matters}{%
\subsubsection{Why This Matters}\label{why-this-matters}}

This transition probability is the \textbf{model} of the environment.
Algorithms that USE this model (Policy Iteration, Value Iteration) are
called \textbf{model-based}. Algorithms that DON'T use it (SARSA, Monte
Carlo, etc.) are called \textbf{model-free} --- they just learn from
trial and error.

\begin{center}\rule{0.5\linewidth}{0.5pt}\end{center}

\hypertarget{the-10-algorithms-explained-one-by-one}{%
\subsection{6. The 10 Algorithms --- Explained One by
One}\label{the-10-algorithms-explained-one-by-one}}

We implemented 10 algorithms across 4 categories. Here's each one
explained simply.

\hypertarget{category-1-dynamic-programming-model-based}{%
\subsubsection{CATEGORY 1: Dynamic Programming
(Model-Based)}\label{category-1-dynamic-programming-model-based}}

These algorithms KNOW the transition probabilities P(s'\textbar s,a).
They use this knowledge to mathematically compute the best policy
without any trial-and-error.

\hypertarget{algorithm-1-policy-iteration-the-winner}{%
\paragraph{Algorithm 1: Policy Iteration ⭐ (THE
WINNER)}\label{algorithm-1-policy-iteration-the-winner}}

\textbf{Analogy}: Imagine you have a study plan. Step 1: Check how well
your current plan works (grade yourself). Step 2: Improve your plan
based on the grades. Repeat until the plan doesn't change anymore.

\textbf{How it works (2 alternating steps)}:

\textbf{Step 1 --- Policy Evaluation}: ``How good is my current
strategy?'' - Take the current policy (strategy) - Calculate V(s) = the
expected total reward starting from state s if you follow this policy -
Uses the formula: V(s) = R(s, π(s)) + γ × Σ P(s'\textbar s, π(s)) ×
V(s') - In English: ``My value = my immediate reward + 0.9× (probability
of each next state × value of that state)'' - Keep updating until values
stop changing

\textbf{Step 2 --- Policy Improvement}: ``Can I do better?'' - For each
state, check all 13 actions - For each action, calculate Q(s,a) =
immediate reward + discounted future value - Pick the action with the
highest Q(s,a) as the new policy - If the policy didn't change → we've
found the optimal policy. STOP!

\textbf{In our project}: Converges in just \textasciitilde4 iterations!
That's because the problem is small enough (244 states) for DP to handle
easily.

\textbf{Why it's the best}: It finds the \textbf{mathematically
guaranteed optimal} policy. No randomness, no approximation, no tuning.

\hypertarget{algorithm-2-value-iteration}{%
\paragraph{Algorithm 2: Value
Iteration}\label{algorithm-2-value-iteration}}

\textbf{Same idea but different approach}: Instead of alternating
between evaluation and improvement, it does both at once in every step.

\textbf{Formula}: V(s) = max\_a {[}R(s,a) + γ × Σ P(s'\textbar s,a) ×
V(s'){]}

In English: ``For each state, find the action that gives the best value,
and set V(s) to that.''

\textbf{Difference from Policy Iteration}: - PI: Fully evaluates the
policy, then improves. Converges in \textasciitilde4 ``big'' iterations.
- VI: Takes tiny steps combining both. Converges in \textasciitilde15
``small'' iterations. - \textbf{Both find the exact same optimal
policy!}

\textbf{Why not chosen as best}: It takes \textasciitilde15 iterations
vs PI's \textasciitilde4. The final answer is the same, but PI is
faster.

\begin{center}\rule{0.5\linewidth}{0.5pt}\end{center}

\hypertarget{category-2-model-free-control}{%
\subsubsection{CATEGORY 2: Model-Free
Control}\label{category-2-model-free-control}}

These algorithms do NOT know P(s'\textbar s,a). They learn by
\textbf{trial and error} --- actually playing the game thousands of
times and learning from experience.

\textbf{Key difference from DP}: Instead of having a ``cheat sheet''
(the transition model), these algorithms are like students who learn by
solving thousands of practice problems.

\hypertarget{algorithm-3-glie-monte-carlo-control}{%
\paragraph{Algorithm 3: GLIE Monte Carlo
Control}\label{algorithm-3-glie-monte-carlo-control}}

\textbf{GLIE} = Greedy in the Limit with Infinite Exploration

\textbf{Analogy}: A student who studies by taking full practice exams,
scores them at the end, and then reviews what went right/wrong.

\textbf{How it works}: 1. Play a FULL episode (from start to diagnosis)
using an ε-greedy policy (usually follow best action, but randomly
explore ε\% of the time) 2. After the episode ends, look back at the
total return G for each (state, action) visited 3. Update: Q(s,a) =
Q(s,a) + (1/N) × (G - Q(s,a)) where N is the visit count 4. Slowly
decrease ε so the AI explores less over time (converges to greedy)

\textbf{Key properties}: - \textbf{Unbiased} --- the return G is the
actual reward the AI received, not an estimate - \textbf{High variance}
--- sometimes lucky, sometimes unlucky episodes give very different
returns - \textbf{Needs 50,000+ episodes} to converge - Updates only
happen AFTER the full episode ends (must wait until diagnosis is made)

\hypertarget{algorithm-4-sarsa-td0-on-policy-control}{%
\paragraph{Algorithm 4: SARSA (TD(0) On-Policy
Control)}\label{algorithm-4-sarsa-td0-on-policy-control}}

\textbf{Name}: S-A-R-S'-A' → State, Action, Reward, next State, next
Action

\textbf{Analogy}: A student who reviews their answers after EVERY
question, not just at the end of the exam.

\textbf{How it works}: 1. In state s, take action a (ε-greedy), get
reward r, end up in state s' 2. In state s', choose next action a'
(ε-greedy) 3. Update immediately: Q(s,a) = Q(s,a) + α × {[}r +
γ×Q(s',a') - Q(s,a){]} 4. Move to s', repeat

\textbf{Key difference vs Monte Carlo}: SARSA doesn't wait for the
episode to end. It updates after EVERY single step. This is called
\textbf{bootstrapping} --- using an estimate (Q(s',a')) instead of the
actual return.

\textbf{Key properties}: - \textbf{Biased} (uses estimate Q(s',a')) but
\textbf{lower variance} than MC - \textbf{On-policy} --- it learns about
the policy it's actually following (including exploration) - Needs α
(learning rate) and ε (exploration rate) --- hyperparameters you must
tune

\hypertarget{algorithm-5-sarsaux3bb-with-eligibility-traces}{%
\paragraph{Algorithm 5: SARSA(λ) --- With Eligibility
Traces}\label{algorithm-5-sarsaux3bb-with-eligibility-traces}}

\textbf{λ (lambda)} blends Monte Carlo and SARSA into one algorithm.

\textbf{Analogy}: Imagine you did something good on question 5, and the
reward came on question 8. MC would only update question 5 after the
exam ends. SARSA would only update question 7 (the step right before the
reward). SARSA(λ) updates MULTIPLE past steps --- how many depends on λ.

\textbf{How it works}: - Keeps a ``trace'' for each (state, action) ---
a record of how recently it was visited - When a reward happens, ALL
recently-visited state-action pairs get updated (proportional to their
trace) - λ=0 → pure SARSA (only update the last step) - λ=1 → pure Monte
Carlo (update all steps equally) - λ=0.8 (our choice) → mostly Monte
Carlo but with some TD bootstrapping

\textbf{Key properties}: - Extra hyperparameter λ to tune - In our
project, got the highest average steps (5.2) --- slightly less efficient
than others

\begin{center}\rule{0.5\linewidth}{0.5pt}\end{center}

\hypertarget{category-3-function-approximation}{%
\subsubsection{CATEGORY 3: Function
Approximation}\label{category-3-function-approximation}}

\textbf{The motivation}: In Categories 1 and 2, we stored Q(s,a) in a
big table with 244×13 = 3,172 entries. This works fine for our small
problem. But for real hospitals with 1000+ diseases and 100+ symptoms?
The table would have billions of entries!

\textbf{The idea}: Instead of a table, use a mathematical FUNCTION to
approximate Q(s,a):

\begin{verbatim}
Q̂(s,a) ≈ φ(s,a)ᵀ × w
\end{verbatim}

Where φ(s,a) is a \textbf{feature vector} (a description of the
state-action pair) and w is a \textbf{weight vector} that gets learned.

\textbf{Our feature vector (215 dimensions)}: - 195 features:
interactions (which symptom × its status × which action) - 13 features:
which action is being taken (one-hot) - 6 features: how many symptoms
are known (binned) - 1 feature: bias term - Total: 195 + 13 + 6 + 1 =
\textbf{215 features}

Instead of learning 3,172 Q-values, we only learn 215 weights. This is
more compact and can generalize across similar states.

\hypertarget{algorithm-6-monte-carlo-with-function-approximation-mc-fa}{%
\paragraph{Algorithm 6: Monte Carlo with Function Approximation
(MC-FA)}\label{algorithm-6-monte-carlo-with-function-approximation-mc-fa}}

Same logic as GLIE MC, but instead of updating a Q-table, it updates the
weight vector w:

\textbf{Update}: w = w + α × (G\_t - φ(s,a)ᵀw) × φ(s,a)

In English: ``Adjust the weights so that the predicted Q-value gets
closer to the actual return.''

\hypertarget{algorithm-7-sarsa-with-function-approximation-sarsa-fa}{%
\paragraph{Algorithm 7: SARSA with Function Approximation
(SARSA-FA)}\label{algorithm-7-sarsa-with-function-approximation-sarsa-fa}}

Same logic as SARSA, but with function approximation:

\textbf{Update}: w = w + α × (r + γ×Q̂(s',a') - Q̂(s,a)) × φ(s,a)

Called ``semi-gradient'' because we only take the gradient of Q̂(s,a) not
Q̂(s',a').

\hypertarget{algorithm-8-lspi-least-squares-policy-iteration}{%
\paragraph{Algorithm 8: LSPI (Least Squares Policy
Iteration)}\label{algorithm-8-lspi-least-squares-policy-iteration}}

A batch method that's quite different:

\begin{enumerate}
\def\labelenumi{\arabic{enumi}.}
\tightlist
\item
  Collect a bunch of experience (10,000 episodes)
\item
  Build matrices A and b from all the data
\item
  Solve w = A⁻¹b (least squares --- like solving a system of equations)
\item
  Extract policy from w
\item
  Collect more data with new policy, repeat
\end{enumerate}

\textbf{Key property}: No learning rate needed! The least squares
solution is exact for the given data. But it needs a lot of data upfront
and is very slow.

\begin{center}\rule{0.5\linewidth}{0.5pt}\end{center}

\hypertarget{category-4-policy-gradient}{%
\subsubsection{CATEGORY 4: Policy
Gradient}\label{category-4-policy-gradient}}

\textbf{The fundamental difference}: All previous algorithms learn
Q-values (how good is each action?) and then derive a policy. Policy
Gradient methods \textbf{directly learn the policy itself} --- no
Q-values at all!

\textbf{The policy is a formula}: π\_θ(a\textbar s) = probability of
taking action a in state s, parameterized by θ.

We use a \textbf{softmax} policy: the probability of each action is
proportional to exp(φ(s,a)ᵀ × θ).

\hypertarget{algorithm-9-reinforce-monte-carlo-policy-gradient}{%
\paragraph{Algorithm 9: REINFORCE (Monte Carlo Policy
Gradient)}\label{algorithm-9-reinforce-monte-carlo-policy-gradient}}

\textbf{How it works}: 1. Play a full episode using the current policy
π\_θ 2. For each (state, action) in the episode, calculate the return
G\_t 3. Update: θ = θ + α × ∇log(π\_θ(a\textbar s)) × G\_t 4.
∇log(π\_θ(a\textbar s)) is the \textbf{score function} --- it tells us
how to change θ to make this action more likely

\textbf{In English}: ``If an action led to a good return (high G\_t),
increase its probability. If it led to a bad return, decrease its
probability.''

\textbf{Key property}: Very high variance because G\_t fluctuates
wildly. Needs 100,000+ episodes.

\hypertarget{algorithm-10-actor-critic}{%
\paragraph{Algorithm 10: Actor-Critic}\label{algorithm-10-actor-critic}}

\textbf{The fix for REINFORCE's variance problem}: Instead of using G\_t
(which is noisy), use a critic to estimate values.

\textbf{Two components}: - \textbf{Actor}: The policy π\_θ(a\textbar s)
--- makes decisions - \textbf{Critic}: A value function V̂(s) ---
evaluates how good the current state is

\textbf{Update}: - \textbf{TD error}: δ = r + γ×V̂(s') - V̂(s) -
\textbf{Critic update}: v = v + α\_critic × δ × ∇V̂(s) - \textbf{Actor
update}: θ = θ + α\_actor × δ × ∇log(π\_θ(a\textbar s))

\textbf{In English}: The critic says ``this state was better/worse than
expected'' (the TD error δ). The actor then adjusts the policy based on
this feedback. Much lower variance than REINFORCE because δ is more
stable than G\_t.

\textbf{Key property}: Two learning rates to tune (one for actor, one
for critic). Sometimes gets 7/8 accuracy instead of 8/8 due to training
instability.

\begin{center}\rule{0.5\linewidth}{0.5pt}\end{center}

\hypertarget{why-policy-iteration-wins-over-everything-else}{%
\subsection{7. Why Policy Iteration Wins Over Everything
Else}\label{why-policy-iteration-wins-over-everything-else}}

\hypertarget{the-short-answer}{%
\subsubsection{The Short Answer}\label{the-short-answer}}

\textbf{Policy Iteration wins because our problem is small enough for it
to work, and when it works, nothing can beat it.}

\hypertarget{the-detailed-reasoning}{%
\subsubsection{The Detailed Reasoning}\label{the-detailed-reasoning}}

\hypertarget{reason-1-we-built-the-transition-model-why-not-use-it}{%
\paragraph{Reason 1: We Built the Transition Model --- Why Not Use
It?}\label{reason-1-we-built-the-transition-model-why-not-use-it}}

We explicitly calculated P(s'\textbar s,a) using the disease patterns
and compatible disease logic. Since we HAVE this model, why would we
throw it away and use model-free methods that ignore it?

\begin{longtable}[]{@{}
  >{\raggedright\arraybackslash}p{(\columnwidth - 4\tabcolsep) * \real{0.31}}
  >{\raggedright\arraybackslash}p{(\columnwidth - 4\tabcolsep) * \real{0.41}}
  >{\raggedright\arraybackslash}p{(\columnwidth - 4\tabcolsep) * \real{0.28}}@{}}
\toprule
Approach & What it does & Like\ldots{} \\
\midrule
\endhead
Policy Iteration & Uses the model to compute exact solution & Having the
answer key and using it \\
SARSA/MC & Ignores the model, learns by trial and error & Throwing away
the answer key and guessing \\
\bottomrule
\end{longtable}

Model-free algorithms are amazing when you DON'T have the model (e.g.,
training a robot to walk --- you can't write down physics equations
easily). But when you DO have the model, using it is always better.

\hypertarget{reason-2-exact-vs-approximate}{%
\paragraph{Reason 2: Exact vs
Approximate}\label{reason-2-exact-vs-approximate}}

\begin{itemize}
\tightlist
\item
  \textbf{Policy Iteration}: Solves the Bellman equations EXACTLY. The
  policy it finds is mathematically proven to be optimal.
\item
  \textbf{SARSA/MC}: Approximate the Q-values through sampling. They
  converge to the optimal given infinite episodes, but in practice with
  finite episodes, there's always some approximation error.
\item
  \textbf{FA methods}: Approximate Q-values with a 215-weight linear
  function. This introduces function approximation error ON TOP OF
  sampling error.
\item
  \textbf{Policy Gradient}: Approximate the policy directly. Even more
  approximation.
\end{itemize}

\textbf{Approximation hierarchy}: PI (exact) \textgreater{} MC/SARSA
(sampling error only) \textgreater{} FA (sampling + approximation error)
\textgreater{} PG (direct policy approximation)

\hypertarget{reason-3-training-speed}{%
\paragraph{Reason 3: Training Speed}\label{reason-3-training-speed}}

\begin{longtable}[]{@{}
  >{\raggedright\arraybackslash}p{(\columnwidth - 4\tabcolsep) * \real{0.38}}
  >{\raggedright\arraybackslash}p{(\columnwidth - 4\tabcolsep) * \real{0.45}}
  >{\raggedright\arraybackslash}p{(\columnwidth - 4\tabcolsep) * \real{0.17}}@{}}
\toprule
Algorithm & Training Time & Why \\
\midrule
\endhead
Policy Iteration & \textasciitilde4 iterations, \textless0.5 seconds &
Directly solves equations \\
Value Iteration & \textasciitilde15 iterations, \textasciitilde2.5
seconds & One-step updates, more iterations \\
GLIE MC & 50,000 episodes, \textasciitilde1.8 seconds & Needs many
episodes, but simple \\
SARSA & 50,000 episodes, \textasciitilde1.8 seconds & Same \\
SARSA(λ) & 50,000 episodes, \textasciitilde2.8 seconds & Traces add
overhead \\
MC-FA & 50,000 episodes, \textasciitilde10 seconds & Feature computation
adds overhead \\
SARSA-FA & 50,000 episodes, \textasciitilde11 seconds & Same \\
LSPI & 10,000 episodes + matrix solve, \textasciitilde63 seconds &
Matrix inversion is expensive \\
REINFORCE & 100,000 episodes, \textasciitilde60 seconds & Very slow
convergence \\
Actor-Critic & 50,000 episodes, \textasciitilde36 seconds & Two networks
to train \\
\bottomrule
\end{longtable}

\hypertarget{reason-4-no-hyperparameters}{%
\paragraph{Reason 4: No
Hyperparameters}\label{reason-4-no-hyperparameters}}

\begin{longtable}[]{@{}ll@{}}
\toprule
Algorithm & Hyperparameters to Tune \\
\midrule
\endhead
\textbf{Policy Iteration} & \textbf{γ only} (and γ=0.9 is standard) \\
Value Iteration & γ \\
GLIE MC & γ, ε\_decay \\
SARSA & γ, α, ε\_decay \\
SARSA(λ) & γ, α, ε\_decay, \textbf{λ} \\
MC-FA & γ, α, ε\_decay \\
SARSA-FA & γ, α, ε\_decay \\
LSPI & γ, ε \\
REINFORCE & γ, α \\
Actor-Critic & γ, α\_actor, α\_critic \\
\bottomrule
\end{longtable}

Hyperparameters are dangerous: wrong values → algorithm doesn't converge
or converges to a bad policy. PI has essentially nothing to tune.

\hypertarget{reason-5-deterministic-decisions}{%
\paragraph{Reason 5: Deterministic
Decisions}\label{reason-5-deterministic-decisions}}

In a medical setting, you want the AI to give the \textbf{same} answer
every time for the same inputs. A patient shouldn't get different
diagnoses based on which random number the AI generated.

\begin{itemize}
\tightlist
\item
  \textbf{PI}: Once trained, the policy is 100\% deterministic. Same
  state → always same action.
\item
  \textbf{Model-free during training}: Use ε-greedy (random
  exploration). Even after training, if ε \textgreater{} 0, there's
  randomness.
\item
  \textbf{Policy Gradient}: The policy is ALWAYS probabilistic
  (softmax). Actions are sampled from a distribution, not chosen
  deterministically.
\end{itemize}

\hypertarget{reason-6-state-space-is-small}{%
\paragraph{Reason 6: State Space Is
Small}\label{reason-6-state-space-is-small}}

244 states × 13 actions = 3,172 state-action pairs. Policy Iteration can
easily handle millions. Our problem is trivially small for it.

FA methods are designed for HUGE state spaces (millions+) where tables
don't fit in memory. Using them here is like using a crane to lift a
coffee cup --- it works, but it's overkill.

\textbf{When WOULD you want FA/PG?} - 50 diseases, 20 symptoms → 3²⁰ ≈
3.5 billion states → table doesn't fit → need FA - Continuous state
spaces (e.g., vital signs as numbers, not binary) → need PG

\hypertarget{summary-table-why-pi-beats-each-algorithm}{%
\subsubsection{Summary Table: Why PI Beats Each
Algorithm}\label{summary-table-why-pi-beats-each-algorithm}}

\begin{longtable}[]{@{}ll@{}}
\toprule
Algorithm & PI is better because\ldots{} \\
\midrule
\endhead
Value Iteration & Same answer, but PI is 3x faster (4 vs 15
iterations) \\
GLIE MC & PI is exact; MC needs 50k episodes and has high variance \\
SARSA & PI is exact; SARSA needs 50k episodes and has α to tune \\
SARSA(λ) & PI is exact; SARSA(λ) has worst avg steps (5.2) and λ to
tune \\
MC-FA & PI is exact; FA adds unnecessary complexity for 244 states \\
SARSA-FA & Same --- FA is overkill for our small state space \\
LSPI & PI is instantaneous; LSPI needs 10k episodes + slow matrix
math \\
REINFORCE & PI is exact; REINFORCE needs 100k episodes, high variance \\
Actor-Critic & PI is exact; A-C sometimes gets 7/8, has 2 learning
rates \\
\bottomrule
\end{longtable}

\begin{center}\rule{0.5\linewidth}{0.5pt}\end{center}

\hypertarget{results-discussion}{%
\subsection{8. Results Discussion}\label{results-discussion}}

\hypertarget{result-1-all-10-algorithms-achieve-100-accuracy}{%
\subsubsection{Result 1: All 10 Algorithms Achieve 100\%
Accuracy}\label{result-1-all-10-algorithms-achieve-100-accuracy}}

Every single algorithm correctly diagnoses all 8 diseases. This
validates: - Our MDP design is sound (diseases are distinguishable) -
The reward structure works (correct=+10, wrong=-5 motivates accuracy) -
All 10 algorithms eventually learn the right thing

\hypertarget{result-2-pi-uses-fewest-steps-4.0-average}{%
\subsubsection{Result 2: PI Uses Fewest Steps (4.0
Average)}\label{result-2-pi-uses-fewest-steps-4.0-average}}

PI asks an average of 4.0 questions per diagnosis. This is essentially
the theoretical minimum needed to distinguish 8 diseases with binary
symptoms (since log₂(8) ≈ 3, and we need slightly more because symptoms
overlap).

\hypertarget{result-3-what-strategy-did-the-ai-learn}{%
\subsubsection{Result 3: What Strategy Did the AI
Learn?}\label{result-3-what-strategy-did-the-ai-learn}}

The AI discovered \textbf{hierarchical elimination} on its own:

\begin{enumerate}
\def\labelenumi{\arabic{enumi}.}
\tightlist
\item
  \textbf{Ask Fever first} --- splits diseases into two groups of 4 each
  (Fever-positive vs Fever-negative)
\item
  \textbf{Ask the most discriminative next symptom} based on which group
  the patient is in
\item
  \textbf{Diagnose after 3-4 questions} when only one disease remains
  compatible
\end{enumerate}

This is exactly what an experienced human doctor would do! The AI
independently rediscovered optimal medical reasoning.

\begin{center}\rule{0.5\linewidth}{0.5pt}\end{center}

\hypertarget{likely-viva-questions-and-answers}{%
\subsection{9. Likely Viva Questions and
Answers}\label{likely-viva-questions-and-answers}}

\textbf{Q: What is the problem you are solving?} A: Building an AI
doctor that diagnoses 8 diseases by asking 5 symptoms in the optimal
order using Reinforcement Learning.

\textbf{Q: What is the state space?} A: 244 states. Each state encodes
the AI's knowledge about 5 symptoms (Unknown/Absent/Present). 3⁵=243
knowledge states + 1 terminal.

\textbf{Q: How do you calculate transition probabilities?} A: We average
over all compatible diseases. If the AI is in a state where
Fever=Present and asks about Cough, we find all diseases that have
Fever. Among those, some have Cough (→ state X) and some don't (→ state
Y). The probability of each next state = fraction of compatible diseases
that lead there.

\textbf{Q: Why did you choose Policy Iteration?} A: Because our state
space is small (244) and we have the complete transition model. PI gives
the mathematically exact optimal policy in just 4 iterations. No other
algorithm can produce a better policy.

\textbf{Q: What if you had 1000 diseases?} A: PI would no longer be
practical because the state space would be too large. We'd use Function
Approximation (SARSA-FA, LSPI) or Policy Gradient methods (Actor-Critic)
that can generalize across states.

\textbf{Q: What is the difference between model-based and model-free?}
A: Model-based (PI, VI) uses P(s'\textbar s,a) --- the transition
probabilities. We explicitly computed these from disease patterns.
Model-free (SARSA, MC) doesn't use P(s'\textbar s,a) --- it learns
purely from actually experiencing episodes.

\textbf{Q: What is the difference between Monte Carlo and SARSA?} A: MC
updates Q-values only after a complete episode ends, using the actual
total return. SARSA updates after every single step, using a
bootstrapped estimate. MC is unbiased but high variance; SARSA is biased
but lower variance.

\textbf{Q: What is Function Approximation? Why do we need it?} A:
Instead of storing Q(s,a) in a table, we approximate it with a function:
Q̂ = features × weights. This allows working with huge state spaces where
a table wouldn't fit in memory. In our 244-state problem, it's overkill,
but we implemented it to show we could.

\textbf{Q: What is the difference between REINFORCE and Actor-Critic?}
A: Both are Policy Gradient methods. REINFORCE uses Monte Carlo returns
(G\_t) → high variance. Actor-Critic uses TD error from a value function
(the critic) → lower variance. Actor-Critic typically converges faster.

\textbf{Q: What does the discount factor γ=0.9 do?} A: It makes future
rewards slightly less valuable than immediate ones. Without it (γ=1),
the AI wouldn't care about how many questions it asks. With γ=0.9,
asking fewer questions (earlier diagnosis) is preferred because later
rewards are discounted.

\textbf{Q: How many questions does the AI need on average?} A: 4.0
questions on average with the optimal Policy Iteration policy. Some
diseases need only 3 questions, some need 4.

\begin{center}\rule{0.5\linewidth}{0.5pt}\end{center}

\emph{This document was prepared for the RL Medical Diagnosis Final
Project by K. Chidwipak (S20230010131).}

\end{document}
